\documentclass{notaaula}

        \titulonota{ENEM: energia e sustentabilidade}
        \creditos{Turma 3A · Ciências da Natureza para o ENEM · Setembro/2026 · Prof. Josué Cavalcante}

        \begin{document}
        \cabecalho

        \begin{sobre}
        Bloco Tecnologia, Sociedade e Intervenção: energia, eletricidade, eletroquímica e impactos socioambientais. Treino integrado de leitura, cálculo e argumentação científica.
        \end{sobre}

        \section{Estratégia de leitura no ENEM}

\begin{copiar}{Do contexto ao modelo}
Antes de calcular, siga quatro perguntas:
\begin{itens}
\item Qual grandeza o item pede e em qual unidade?
\item Quais dados do texto, gráfico ou tabela são realmente necessários?
\item Qual transformação de energia descreve o sistema?
\item A resposta é compatível com a ordem de grandeza e com o impacto ambiental discutido?
\end{itens}
O ENEM costuma avaliar seleção de dados e interpretação mais do que substituição direta em fórmulas.
\end{copiar}

\section{Energia, potência e eficiência}

\begin{copiar}{Núcleo quantitativo}
\[
  E=P\Delta t,
  \qquad P=Ui,
  \qquad \eta=\frac{E_{\mathrm{útil}}}{E_{\mathrm{total}}}.
\]
Use kW e h para obter kWh. Compare alternativas energéticas por eficiência, disponibilidade,
regularidade, custo, emissões, ocupação do território e resíduos. Nenhuma fonte é totalmente sem impacto.
\end{copiar}

\begin{exemplo}[Item comentado]
Uma escola troca 100 lâmpadas de $60\un{W}$ por modelos de $10\un{W}$, usados $5\un{h/dia}$ em
20 dias letivos. A economia mensal é
\[
  \Delta E=100(\dec{0,060}-\dec{0,010})\cdot100
  =\resultado{500\un{kWh}}.
\]
A melhor conclusão combina o cálculo com menor demanda elétrica, sem afirmar impacto ambiental nulo.
\end{exemplo}

\section{Pilhas, eletrólise e materiais}

\begin{copiar}{Conversões eletroquímicas}
Em pilhas, reação espontânea converte energia química em elétrica. Na eletrólise, uma fonte elétrica
força reação não espontânea. Oxidação perde elétrons; redução ganha elétrons.
\begin{itens}
\item Pilhas e baterias exigem coleta e descarte adequados por causa dos materiais reativos.
\item Recarga não elimina impactos: origem da eletricidade, vida útil e reciclagem importam.
\item Comparações devem usar a mesma função e o ciclo de vida, não apenas a etapa de uso.
\end{itens}
\end{copiar}

\section{Energia em sistemas vivos}

\begin{copiar}{Integração com Biologia e Química}
Fotossíntese armazena energia luminosa em ligações químicas; respiração celular transfere parte dessa
energia para ATP. Cadeias alimentares transferem energia com perdas na forma de calor a cada nível
trófico. A 1ª lei da Termodinâmica organiza o balanço; a 2ª lei ajuda a compreender a degradação da
energia e a impossibilidade de rendimento de $100\%$ em processos reais.
\end{copiar}

\begin{dica}
Desconfie de alternativa absoluta: ``impacto zero'', ``todo'', ``nunca'' ou ``100\% eficiente'' costuma
contrariar os balanços de matéria e energia do próprio enunciado.
\end{dica}

\begin{exercicios}[Mini-simulado ENEM]
\begin{questoes}
\item Um aparelho de $2\un{kW}$ opera por 30 minutos. Seu consumo é
  \begin{alternativas}\item $\dec{0,5}\un{kWh}$\item $1\un{kWh}$\item $2\un{kWh}$\item $4\un{kWh}$\end{alternativas}
\item Uma máquina recebe $500\un{J}$ e entrega $350\un{J}$ úteis. O rendimento é
  \begin{alternativas}\item $30\%$\item $50\%$\item $70\%$\item $85\%$\end{alternativas}
\item Na pilha em funcionamento, a conversão principal é
  \begin{alternativas}\item elétrica em química\item química em elétrica\item térmica em nuclear\item luminosa em mecânica\end{alternativas}
\item A afirmação cientificamente mais adequada sobre energia solar é
  \begin{alternativas}\item não produz impacto em nenhuma etapa\item dispensa armazenamento sempre\item reduz emissões na operação, mas materiais e área devem ser avaliados\item tem rendimento de $100\%$\end{alternativas}
\item Em cadeia alimentar, a energia disponível diminui nos níveis superiores principalmente porque
  \begin{alternativas}\item energia é destruída\item parte se dissipa como calor em processos metabólicos\item matéria não circula\item produtores não respiram\end{alternativas}
\item Trocar equipamento por outro de mesma função e menor potência reduz consumo quando
  \begin{alternativas}\item o tempo de uso não aumenta o bastante para anular a economia\item a tensão é sempre zero\item a tarifa deixa de existir\item o rendimento supera $100\%$\end{alternativas}
\item Explique por que análise de ciclo de vida é melhor que observar apenas o uso de uma bateria.
\item Uma escola economiza $120\un{kWh}$ a tarifa de $R\$\,\dec{0,85}/\un{kWh}$. Calcule a economia financeira.
\end{questoes}
\gabarito{1-b; 2-c; 3-b; 4-c; 5-b; 6-a; 8) $R\$\,102{,}00$.}
\end{exercicios}


\end{document}
